\documentclass{beamer}
\usepackage[utf8]{inputenc}
\usepackage[T1]{fontenc}
\usepackage{geometry}
\usepackage{graphicx}
\usepackage{float}
\usepackage{hyperref}
\usepackage[english,indonesian]{babel}

\title{Optimasi Pemrosesan Template LaTeX Menggunakan Pendekatan Berbasis AI}
\author{Academic Support Agent}
\date{\today}

\begin{document}
\maketitle

\begin{abstract}

\end{abstract}

\chapter{PENDAHULUAN}


\chapter{TINJAUAN PUSTAKA}


% Original content (commented out):
% \documentclass[10pt,aspectration=169]{beamer}
%% Basic weizman color 34B0E5 - RGB[52,176,229] RGB%[20.4,69,89.8]
%% DARK  0483B9, {RGB}{4,131,185} RGB: 4–131–185
\usetheme[
%%% option passed to the outer theme %    progressstyle=fixedCircCnt,   % fixedCircCnt, movingCircCnt (moving is deault) % If you want to change the colors of the various elements in the theme, edit and uncomment the following lines
  ]{Feather}
  \definecolor{Wblue}{RGB}{4,131,185}    %{199,97,107} 
% Change the bar colors:
\setbeamercolor{Feather}{fg=Wblue!20,bg=Wblue}

% Change the color of the structural elements:
\setbeamercolor{structure}{fg=Wblue}
% Change the frame title text color:
%\setbeamercolor{Feather}{fg=Praga}
%\setbeamercolor{frametitle}{fg=blue}

% Change the normal text color background:
%\setbeamercolor{normal text}{fg=black,bg=gray!10}

%-------------------------------------------------------
% INCLUDE PACKAGES
%-------------------------------------------------------
\usepackage[version=3]{mhchem}
\usepackage[utf8]{inputenc}
\usepackage[english]{babel}
\usepackage[T1]{fontenc}
\usepackage{helvet}
\usepackage{mathtools}
\usefonttheme{professionalfonts}
\usepackage{lmodern}
\usepackage{bm}
\usepackage{tikz}
%\usepackage{soul} %strikethrough
\usepackage{siunitx}
\usepackage{times}
\usepackage{tikz}
\usepackage{amsmath}
\usepackage{verbatim}
\usetikzlibrary{arrows,shapes}
\usepackage{amsmath}
\usepackage{booktabs}
\usepackage{amsfonts}
\usepackage{amssymb}
\usepackage{graphicx}
\usepackage{cancel}
\usepackage{multirow}
\DeclareTextFontCommand{\emph}{\bfseries}
\setbeamertemplate{bibliography item}{\insertbiblabel}
%citations.bib
\usepackage[backend=bibtex,
style=numeric,
maxnames=2,
natbib=true,
abbreviate=false,
sorting=none,
giveninits=true,
autocite=superscript
]{biblatex}
\addbibresource{main.bib}
\renewcommand*{\bibfont}{\tiny}
%\Atbeginbibliography{\tiny}
%-------------------------------------------------------
% DEFFINING AND REDEFINING COMMANDS
%-------------------------------------------------------
\DeclareBibliographyDriver{article}{\printnames{author}   \newunit     \printfield{year}     \newunit     \printfield{journaltitle}     \newunit     \printfield{volume}     \newunit     \printfield{pages}     \finentry}
\newcommand{\chref}[2]{
  \href{#1}{{\usebeamercolor[bg]{Feather}#2}}
}

% INFORMATION IN THE TITLE PAGE

\title[SHORT TITLE] % [] is optional - is placed on the bottom of the sidebar on every slide
{ % is placed on the title page
      \textbf{TITLE}
}

%\subtitle[The Feather Beamer Theme]
\institute[]
{
    University Name, Town, State\\
}

\date{\today}

\begin{document}

{\1% % this is the name of the PDF file for the background
\begin{frame}[plain,noframenumbering] % the plain option removes the header from the title page, noframenumbering removes the numbering of this frame only
\begin{center}

 \textbf{\Large \\[2em]Here we put a very nice title\\[1em]}
 
 \end{center}

    \begin{center} \line(1,0){250} \end{center}

\textbf{Author One*, Author Two, Final Author }\\


	\begin{center}
 \textbf{LOCATION AND DATE} \\
 \vspace{0.05\textheight}
\includegraphics[width=0.98\textwidth]{./im/001_1pagina_sanzio.jpg}

	\end{center}

\begin{tiny}
Raffaello Sanzio - La scuola di Atene
\end{tiny}

\end{frame}}

\begin{frame}\frametitle{Outline}
Remember that the outline is not:
\begin{itemize}
    \item Intro
    \item Theory
    \item Data
    \item Conclusion
\end{itemize}
That is just the logical structure of ANY presentation.
A proper Outline is:
\begin{itemize}
    \item We want to save the world
    \item SuperHeroes we considered 
    \item IronMan is the best solution
    \item Strenght is the best attribute of IronMan
    \item We saved the world thanks to IronMan
\end{itemize}
\end{frame}

\begin{frame}{Aim of the project}
ENUMERATE
\begin{enumerate}
    \item One
    \item Two 
\end{enumerate}
\begin{small}
\begin{center}

I like to split data and figure in columns: \\
\begin{columns}

\column{0.33\textwidth}
First column
\column{0.33\textwidth}
Maybe a figure here
\column{0.33\textwidth}
\begin{itemize}
    \item[] Usually
    \item[] 3 columns
    \item[] are one
    \item[] too many
\end{itemize}
\end{columns}
\end{center}
\end{small}

\end{frame}
% Hook is: looking different aspect, we start from 
%\begin{frame}{Methodology}
%I will present calculations results using ab-initio, Density Functional Theory. \\
%
%\end{frame}

\begin{frame}{Biblio}
    
\begin{small}
Fig from\footfullcite{ambacher2000two}. P Data from\footfullcite{bernardini1997spontaneous}
\end{small}
\addtocounter{footnote}{-1}
\end{frame}
%%%%%%%%%%%%
\end{document}


\end{document}